% ===========================================================================
%  Chapter 04 homework — Field extensions
%  Garling, A Course in Galois Theory
%
%  DIVISION OF LABOUR (see AI_INSTRUCTIONS.md §7):
%    The assistant transcribes each assigned problem statement faithfully and
%    emits an empty, marked solution region beneath it.
%    The user types the mathematics inside those regions. No assistant writes
%    inside a SOLUTION region in normal mode.
% ===========================================================================
\documentclass[11pt]{article}
\usepackage{coursemacros}

\title{Chapter 04 Homework --- Field extensions}
\author{Mikey Ferguson}
\date{\today}

\begin{document}
\maketitle

% ---------------------------------------------------------------------------
% Assigned set: all eleven exercises of chapter 4, 4.1 through 4.11.
% Printed here in the book's order. The sitting works them in a different
% order (4.1, 4.4, 4.6, 4.5, then 4.2, 4.3, then 4.7-4.11).
% Statements are all transcribed; solution regions stay empty until the
% student's answer is agreed correct.
% ---------------------------------------------------------------------------

\begin{problem}{04.1}
  Suppose that $\degree{L}{K}$ is a prime number. What fields are there
  intermediate between $L$ and $K$?
\end{problem}
% ===== SOLUTION 04.1 =====
% Transcribed from the student's slate, 2026-09-15 (pages 01-03, turns leading
% to cards 0003-0006). Typeset exactly as written: same steps, same order,
% nothing added.
% One notation change only, agreed on card 0004: the scalar with
% k(l_1) = 1 was called k_1 on the page, the same name as the general
% coefficient. It is called k_0 here, as on card 0003.
% Not covered by what they wrote, and flagged to them on card 0005:
%   - the closing line that K =/= L is given verbally, not on the page;
%     the page stops at the list.
\begin{solution}
  \textbf{Lemma.} If $\degree{L}{K} = 1$ then $L = K$.

  If $\degree{L}{K} = 1$, that means you only need one vector of $L$ to span
  all of $L$ if you are given any coefficients in $K$. Suppose $\set{\ell_{1}}$
  spans $L$. Then
  \[
    \forall\, \ell \in L,\ \exists\, k \in K \text{ s.t. } k\ell_{1} = \ell .
  \]
  So $k_{0}\ell_{1} = 1$ for some $k_{0} \in K$, and surely $k_{0} \neq 0$,
  since that forces $k_{0}\ell_{1} = 0$. Hence $k_{0} = \ell_{1}^{-1}$, so
  $\ell_{1}^{-1} \in K$, so $\ell_{1} \in K$, since $K$ is a field.

  Let $\ell \in L$; show $\ell \in K$. There is $k_{1}$ with
  $k_{1}\ell_{1} = \ell$. Both $k_{1}, \ell_{1} \in K$, so $\ell \in K$.

  \medskip
  \textbf{The exercise.} Let $\degree{L}{K} = p$ be prime and let $M$ be
  intermediate, $K \subseteq M \subseteq L$. The tower law gives
  \[
    \degree{L}{K} = \degree{L}{M}\,\degree{M}{K},
  \]
  and the only factorisations of $p$ into positive degrees are
  \[
    \begin{array}{ccc}
      p & 1 & p \\
      p & p & 1
    \end{array}
  \]
  that is, $\degree{L}{M} = 1$ with $\degree{M}{K} = p$, or $\degree{L}{M} = p$
  with $\degree{M}{K} = 1$. By the lemma the first gives $M = L$ and the
  second gives $M = K$. Both occur:
  \[
    \degree{L}{K} = \degree{L}{L}\,\degree{L}{K}
    \qquad\text{and}\qquad
    \degree{L}{K} = \degree{L}{K}\,\degree{K}{K}.
  \]

  Which fields are intermediate? Just $L$ and $K$.
\end{solution}
% ===== END SOLUTION 04.1 =====

\begin{problem}{04.2}
  Suppose that $\ext{L}{K}$ and that $K_{1}$ and $K_{2}$ are two intermediate
  fields such that $L = \adjoin{K}{K_{1}, K_{2}}$. Show that
  $\degree{L}{K} \leq \degree{K_{1}}{K}\,\degree{K_{2}}{K}$.
\end{problem}
% ===== SOLUTION 04.2 =====
% TODO(mferguson): your work goes here.
% ===== END SOLUTION 04.2 =====

\begin{problem}{04.3}
  Suppose that $\ext{\adjoin{K}{\alpha}}{K}$ is a finite simple extension. For
  each $\beta$ in $\adjoin{K}{\alpha}$, let $T_{\alpha}(\beta) = \alpha\beta$.
  $T_{\alpha}$ is a linear mapping of $\adjoin{K}{\alpha}$ (considered as a
  vector space over $K$) into itself. Show that $\det(xI - T_{\alpha})$ is the
  minimal polynomial of $\alpha$ over $K$.
\end{problem}
% ===== SOLUTION 04.3 =====
% TODO(mferguson): your work goes here.
% ===== END SOLUTION 04.3 =====

\begin{problem}{04.4}
  Show that $x^{3} + 3x + 1$ is irreducible in $\polyring{\QQ}{x}$. Suppose
  that $\alpha$ is a root of $x^{3} + 3x + 1$ in $\CC$. Express $\alpha^{-1}$
  and $(1 + \alpha)^{-1}$ as linear combinations, with rational coefficients,
  of $1$, $\alpha$ and $\alpha^{2}$.
\end{problem}
% ===== SOLUTION 04.4 =====
% Transcribed from the student's slate, 2026-09-15 (pages 04-07, turns leading
% to cards 0007-0010). Typeset exactly as written: same steps, same order,
% nothing added.
% Not covered by what they wrote, and flagged to them on cards 0007 and 0009:
%   - why the only candidate rational roots are +1 and -1 (monic, constant 1)
%     is used but not stated;
%   - "no rational root, hence irreducible" needs the degree-three hypothesis
%     to be sound; the page asserts the implication without it.
\begin{solution}
  \textbf{Irreducibility.} $f(x) = x^{3} + 3x + 1 \in \polyring{\QQ}{x}$, so
  the possible rational roots are $\set{1, -1}$.
  \[
    f(-1) = -1 - 3 + 1 \neq 0, \qquad f(1) = 1 + 3 + 1 \neq 0 .
  \]
  So $f(x)$ has no rational roots and is hence irreducible in
  $\polyring{\QQ}{x}$.

  \medskip
  \textbf{Inverting $\alpha$.} $\alpha$ is a root, so
  $f(\alpha) = \alpha^{3} + 3\alpha + 1 = 0$. Then
  \[
    \alpha^{3} + 3\alpha + 1 = 0
    \iff
    -1 = \alpha(\alpha^{2} + 3),
    \quad\text{so}\quad
    \alpha^{-1} = \underbrace{-1}_{\in\QQ}\cdot\,\alpha^{2}
                  \;\underbrace{-\;3}_{\in\QQ}\cdot\,1 .
  \]

  \medskip
  \textbf{Inverting $1 + \alpha$.} Set $\beta = 1 + \alpha$, so
  $\alpha = \beta - 1$. From $\alpha^{3} + 3\alpha + 1 = 0$,
  \begin{align*}
    (\beta - 1)^{3} + 3(\beta - 1) + 1 &= 0 \\
    \iff (\beta - 1)(\beta^{2} - 2\beta + 1) + 3\beta - 3 + 1 &= 0 \\
    \iff \beta^{3} - 2\beta^{2} + \beta - \beta^{2} + 2\beta - 1
         + 3\beta - 3 + 1 &= 0 \\
    \iff \beta^{3} - 3\beta^{2} + 6\beta - 3 &= 0 \\
    \iff 1 &= \tfrac{1}{3}\beta^{3} - \beta^{2} + 2\beta \\
    \iff 1 &= \beta\bigl(\tfrac{1}{3}\beta^{2} - \beta + 2\bigr),
  \end{align*}
  so
  \begin{align*}
    \beta^{-1} = (\alpha + 1)^{-1}
      &= \tfrac{1}{3}(\alpha + 1)^{2} - (\alpha + 1) + 2 \\
      &= \tfrac{1}{3}(\alpha^{2} + 2\alpha + 1) - \alpha + 1 \\
      &= \underbrace{\tfrac{1}{3}}_{\in\QQ}\,\alpha^{2}
         \;\underbrace{-\;\tfrac{1}{3}}_{\in\QQ}\,\alpha
         \;+\;\underbrace{\tfrac{4}{3}}_{\in\QQ}\cdot 1
       = (\alpha + 1)^{-1}.
  \end{align*}
\end{solution}
% ===== END SOLUTION 04.4 =====

\begin{problem}{04.5}
  Suppose that $\ext{\ext{\adjoin{L}{\alpha}}{L}}{K}$ and that
  $\degree{\adjoin{L}{\alpha}}{L}$ and $\degree{L}{K}$ are relatively prime.
  Show that the minimal polynomial of $\alpha$ over $L$ has its coefficients
  in $K$.
\end{problem}

\begin{remark}
  As printed, 4.5 is false. The tower $K = \QQ$, $L = \QQ(2^{1/3})$,
  $\alpha = 2^{1/6}$ has $\degree{\adjoin{L}{\alpha}}{L} = 2$ and
  $\degree{L}{K} = 3$, which are coprime, while the minimal polynomial of
  $\alpha$ over $L$ is $x^{2} - 2^{1/3}$, whose constant term is not rational.
  The repaired hypothesis is that
  $\degree{\adjoin{K}{\alpha}}{K}$ and $\degree{L}{K}$ are relatively prime.
  Established in the sitting of 2026-09-15, card 0016.
\end{remark}
% ===== SOLUTION 04.5 =====
% TODO(mferguson): your work goes here.
% In progress as of 2026-09-15 (slate page 11, answers t0050-r1..r15).
% ===== END SOLUTION 04.5 =====

\begin{problem}{04.6}
  Suppose that $\degree{L}{K}$ is a prime number. Show that $\ext{L}{K}$ is
  simple.
\end{problem}
% ===== SOLUTION 04.6 =====
% Transcribed from the student's slate, 2026-09-15 (page 10, turn leading to
% card 0013). Typeset exactly as written: same steps, same order, nothing
% added.
% Not covered by what they wrote, and supplied on card 0013:
%   - why K(alpha) =/= K: alpha lies in K(alpha) by definition and alpha is
%     not in K;
%   - why an l outside K exists at all: p is prime so p >= 2, so L =/= K.
\begin{solution}
  $\degree{L}{K} = p$. $L$ and $K$ are the only intermediate fields
  (Exercise 4.1).

  Note $\adjoin{K}{\alpha}$ is by definition the smallest subfield containing
  $\alpha \in L$ and all $k \in K$.

  Pick $\alpha = \ell$ where $\ell \notin K$, which exists since
  $\degree{L}{K} > 1$. Note $K \subsetneq \adjoin{K}{\alpha}$.

  So $\adjoin{K}{\alpha}$ must be $L$, the only other intermediate subfield
  between $K$ and $L$.
\end{solution}
% ===== END SOLUTION 04.6 =====

\begin{problem}{04.7}
  Show that if $\ext{L}{K}$ is algebraic and $K$ is countable then $L$ is
  countable. Show that there exist real numbers which are transcendental over
  the rationals.
\end{problem}
% ===== SOLUTION 04.7 =====
% TODO(mferguson): your work goes here.
% ===== END SOLUTION 04.7 =====

\begin{problem}{04.8}
  Suppose that $\ext{L}{K}$ is an extension, that $\alpha$ is an element of
  $L$ which is transcendental over $K$, and that $f$ is a non-constant element
  of $\polyring{K}{x}$. Show that $f(\alpha)$ is transcendental over $K$. Show
  that, if $\beta$ is an element of $L$ which satisfies $f(\beta) = \alpha$,
  then $\beta$ is transcendental over $K$.
\end{problem}
% ===== SOLUTION 04.8 =====
% TODO(mferguson): your work goes here.
% ===== END SOLUTION 04.8 =====

\begin{problem}{04.9}
  Suppose that $a$ and $b$ are complex numbers which are transcendental over
  $\QQ$. Is $ab$ transcendental over $\QQ$?
\end{problem}
% ===== SOLUTION 04.9 =====
% TODO(mferguson): your work goes here.
% ===== END SOLUTION 04.9 =====

\begin{problem}{04.10}
  Suppose that $\ext{\adjoin{K}{\alpha, \beta}}{K}$ is an extension, that
  $\alpha$ is algebraic over $K$, but not in $K$, and that $\beta$ is
  transcendental over $K$. Show that $\ext{\adjoin{K}{\alpha, \beta}}{K}$ is
  not simple.
\end{problem}
% ===== SOLUTION 04.10 =====
% TODO(mferguson): your work goes here.
% ===== END SOLUTION 04.10 =====

\begin{problem}{04.11}
  Show that the condition that $\ext{L}{K}$ is algebraic cannot be dropped
  from Theorem 4.8.
\end{problem}
% ===== SOLUTION 04.11 =====
% TODO(mferguson): your work goes here.
% ===== END SOLUTION 04.11 =====

\end{document}
